%% ============================================================================
\section{Landscape-Scale Inference}
\label{ch:discovery}
%% ============================================================================

The experiments reported so far measure how efficiently a model can be
\emph{trained}. This section describes how the trained models were
\emph{applied}: a single inference campaign over 55{,}008~km$^2$ of the southern
Mesopotamian floodplain, run with all three architectures, producing
georeferenced raster layers. The account below is confined to what was
implemented and how the data was processed.

\subsection{Survey area and tiling}
\label{sec:disc-area}

The survey area is defined by two vector layers: a polygon delimiting the
region of interest, and an exclusion layer covering already-catalogued sites. A
1~km grid was laid over the bounding box of the survey polygon in EPSG:32638
(WGS~84 / UTM zone~38N), giving 322 rows by 395 columns. A cell was retained
only if at least $1\%$ of its area fell inside the survey polygon and it did not
intersect the exclusion layer, leaving 57{,}363 eligible cells.

Each retained cell was rendered from Bing satellite imagery at
$1024\times1024$~px, giving a ground sample distance of $0.977$~m per pixel.
This matches the training corpus: the Bing~1k tiles described in
Section~\ref{ch:materials} are also 1~km rendered at 1024~px. Rendering halted
after 55{,}008 of the 57{,}363 eligible cells; because cells are numbered in
row-major order from the southern edge, the 2{,}355 unrendered cells form the
northernmost strip rather than a scattered sample.

\subsubsection{Manifest reconstruction}
The rendering process recorded the imagery but not the geographic extent of each
tile, leaving 55{,}008 files named \code{tile\_000001.jpg} onwards with no
information about where any of them was. The mapping was recovered by replaying
the renderer's eligibility rule over the same two vector layers in the same
order, which regenerates the tile sequence deterministically. The replay
reproduces exactly 55{,}008 consecutive identifiers with no gaps and no surplus
files, and records the SHA-256 digest of both input layers alongside the grid
origin and dimensions, so the geometry can be re-derived and checked.

\subsection{Inference procedure}
\label{sec:disc-procedure}

Each $1024\times1024$ tile was divided into four $512\times512$ windows at a
stride of 512, so the windows tile the image exactly: every pixel is predicted
once, with no overlap to blend and no gaps to fill. Each window was resampled
bilinearly to $256\times256$, passed through the network, and the resulting
probability map resampled back to $512\times512$ before being written into its
position in the tile.

This reproduces the training-time preprocessing. Training applies a
$512\times512$ random crop followed by a resize to $256\times256$; inference
applies a $512\times512$ deterministic crop followed by the same resize, so the
network receives imagery at the same effective ground sample distance it was
trained on. A sliding window at native resolution would instead have presented
it with imagery at twice that scale.

The weights used are the full-protocol reproduction checkpoints of
Table~\ref{tab:repro-full}, not the checkpoints from the \pads{} staging
experiments, which were trained under an abbreviated protocol intended for
throughput measurement.

Per-tile probability maps were written to node-accessible scratch at half
precision. This halves the intermediate footprint from 692~GB to 346~GB, at a
maximum quantisation error of $2.4\times10^{-4}$ -- three orders of magnitude
below the $0.5$ decision threshold used subsequently.

\subsection{Mosaicking and outputs}
\label{sec:disc-outputs}

The per-tile predictions were assembled into four georeferenced rasters
covering a $387{,}072\times296{,}960$ pixel extent ($378\times290$~km) in
EPSG:32638. Three are single-band float32 rasters holding each model's
per-pixel probability. The fourth composites them: each model's probability is
thresholded at $0.5$ and the three binary decisions encoded as a three-bit
value, giving eight classes -- no detection, each model alone, each of the three
pairs, and agreement of all three -- rendered as a four-band RGBA image.

Roughly half of the bounding box lies outside the survey polygon. The rasters
are written sparsely with deflate compression, so unsurveyed area costs no
storage, and overview pyramids were built with nearest-neighbour resampling,
which preserves the categorical class colours that averaging would blend.

Retaining the continuous probabilities alongside the thresholded composite is
what allows the $0.5$ threshold to be revisited without repeating inference.

%% ---------------------------------------------------------------------------
%% PLACEHOLDER FIGURE 1 -- regional view of the agreement layer.
%% To insert the real figure: save the screenshot as
%%   thesis/figures/discovery-overview.png
%% then delete the \fbox{...} block below and uncomment the \includegraphics.
%% ---------------------------------------------------------------------------
\begin{figure}[!ht]
\centering
%\includegraphics[width=0.95\linewidth]{discovery-overview.png}
\fbox{\parbox[c][6cm][c]{0.95\linewidth}{\centering\textit{Placeholder --
regional view of the model-agreement layer. Replace with
\code{discovery-overview.png}.}}}
\caption[Model-agreement layer over the survey area]{The model-agreement layer
over the surveyed portion of the floodplain. Colour encodes which of the three
architectures exceeded the $0.5$ threshold at each pixel; unanimous agreement is
shown in white and unsurveyed ground is transparent.}
\label{fig:disc-overview}
\end{figure}

%% ---------------------------------------------------------------------------
%% PLACEHOLDER FIGURE 2 -- individual detections at tile scale.
%% Suggested content: two or three 1 km tiles side by side, each showing the
%% Bing imagery with the agreement overlay on top. Choose at least one tile
%% where all three models agree, and one where only a single model fired.
%% Save as thesis/figures/discovery-candidates.png, then delete the \fbox{...}
%% block and uncomment the \includegraphics.
%% ---------------------------------------------------------------------------
\begin{figure}[!ht]
\centering
%\includegraphics[width=\linewidth]{discovery-candidates.png}
\fbox{\parbox[c][6cm][c]{0.95\linewidth}{\centering\textit{Placeholder --
detections at tile scale, imagery with agreement overlay. Replace with
\code{discovery-candidates.png}.}}}
\caption[Detections at tile scale]{Detections at the scale of individual 1~km
tiles, with the agreement overlay drawn over the underlying satellite imagery.
Tiles on which all three architectures agree are contrasted with tiles on which
only one responded.}
\label{fig:disc-candidates}
\end{figure}

%% ---------------------------------------------------------------------------
%% PLACEHOLDER FIGURE 3 -- optional. Per-model probability for one tile.
%% Suggested content: one tile, four panels -- imagery, then the DeepLabV3+,
%% MA-Net and SegFormer probability maps as continuous greyscale or heat maps.
%% Save as thesis/figures/discovery-permodel.png. Delete this whole figure
%% environment if it is not used.
%% ---------------------------------------------------------------------------
\begin{figure}[!ht]
\centering
%\includegraphics[width=\linewidth]{discovery-permodel.png}
\fbox{\parbox[c][5cm][c]{0.95\linewidth}{\centering\textit{Placeholder --
per-model probability maps for a single tile. Replace with
\code{discovery-permodel.png}, or delete this figure.}}}
\caption[Per-model probability for a single tile]{Continuous probability output
of each architecture for a single tile, before thresholding. These differences
are what the agreement classes summarise and what the thresholded composite
discards.}
\label{fig:disc-permodel}
\end{figure}

\subsection{Computational cost}
\label{sec:disc-cost}

The campaign required 660{,}096 forward passes (55{,}008 tiles $\times$ four
windows $\times$ three models). Inference occupied approximately $2.4$~hours on
a single V100. Mosaicking took a further $6$~hours and $19$~minutes on a CPU
node -- more than twice the cost of the inference it was assembling.

That imbalance is the phenomenon this thesis documents for training, appearing
again at deployment. Measured on the inference workload, GPU-busy time was
$22$--$45\%$ depending on architecture; the remainder was spent reading and
decoding JPEG tiles from the shared filesystem, one at a time, in the same
thread as the computation. The mosaicking stage is entirely filesystem-bound. A
faster GPU would have changed the total campaign time very little.

The storage requirement is likewise dominated not by the imagery ($49$~GB) but
by the intermediate predictions, since every model writes a floating-point
probability for every pixel it examines. Inference over a large area is a
storage-provisioning problem before it is a compute problem.

Nothing produced by this campaign has been verified against the ground; the
limits stated in Section~\ref{sec:arch-verification} apply unchanged.
